% s-127-p127-m005; original PDF 127, 128, 129
\Aphorism{s-127-p127-m005}{174}
Die Liebe hat sich vom Fortpflanzungsgeschäft ebensoweit entfernt und ist ihren eigenen Subtilitäten nachgegangen, wie der menschliche Intellect vom Geschäft, Futter zu suchen und Beute zu machen; jedes Organ ist, über den ursprünglichen Naturzweck hinaus, eines selbständigen Daseins, mit all seinen Leiden und Freuden, fähig. Der Mensch ist das supranaturalistische Thier, und in der Liebe ist er es in verhängnissvoller Weise. Er findet am dumpfen, stumpfen Geschlechtsgenuss kein Genügen mehr. Er stellt specielle, endlich gar individuelle Bedingungen. Er wird zart, raffinirt, wählerisch. Ein Zusammenklang von Einzelheiten entzückt ihn; die geringste Störung dieses Ensembles, und er erkaltet, ist dégoûtirt. Es muss ein schwarzes Auge sein, schwarz-blau, nicht schwarzbraun. Das Auge muss gross sein, nicht zu gross, und von einer Braue mit bestimmtem Krümmungsradius überwölbt, dem ein gewisser Krümmungsradius der Mundwölbung entsprechen muss. Die Haut von der Farbe des Elfenbeins, aber ohne Anämie zu verrathen. Die Nase — nun beginnen die Hauptschwierigkeiten; und immer feiner, immer haarschärfer werden die Grenzen, die um kein Zuviel, kein Zuwenig verfehlt werden dürfen — nicht die Grenzen der objectiven Schönheit, wie sich von selbst versteht, sondern der individuellen Erotik, das Liniennetz, in dem gerade dieser Mann gefangen wird, der Specialfall von Constitution, für den drüben, beim andern Geschlecht, der correspondirende Specialfall bereitsteht. Ist diese Correspondenz gegenseitig, nun, so zweifle man nicht, dass hier zwei Organismen einander ausgesuchte und unersetzliche Freuden bereiten, und dass Jeder den Andern für das Zugespitzte und Singuläre seiner Ansprüche durch ebenso zugespitzte und singuläre Befriedigung entschädigt; diese Zwei dürfen schon ein Leben lang auf einander warten. Aber meistens ist die Correspondenz einseitig, und die erotischen Genies gerathen mit ihrer Liebe an die erotischen Dutzendmenschen. Ein Weib, das allen Sublimitäten eines erotisch sehr differenzirten Mannes gerecht wird, physiologisch gerecht wird, kann psychisch genommen ein blosses Fortpflanzungsgeschöpf sein und nach einem undifferenzirten Zuchtstier von Mann Verlangen tragen. Dies ist vielleicht die Regel; denn, unter uns gefragt, giebt es erotisch begabte Weiber? Solche, die in der Liebe eine Kunst, einen seltenen Feiertag der Seele empfinden und mit sensitiver Feinheit an ihre Ausgestaltung gehen? Das Weib, das sich noch die Naivität und Dumpfheit der Gattung bewahrt hat, steht dem bewusst erotischen, wählenden, schmeckenden Manne gegenüber biologisch höher, aber ästhetisch tiefer: es wird diesen Mann nie erhören, — aber oft genug fasciniren, und dann ist das Verhängniss da. Das erotische Genie pflanzt sich so wenig fort, wie die anderen Arten der Genialität.
\par

